\section{The arithmetic operator and its normalization}
\label{sec:normalization}
This appendix fixes the exact external input used in
Proposition~\ref{prop:LPS}, then proves the remaining representation
and normalization steps. Write $\mathbf i,\mathbf j,\mathbf k$ for
quaternion units and
\[
 \mathscr Q_5=\{1\pm2\mathbf i,1\pm2\mathbf j,1\pm2\mathbf k\},
 \qquad r_q(z)=qzq^{-1}\quad(z\in\operatorname{Im}\mathbb H).
\]
The required external LPS spherical Hecke inequality, specialized
to this representative convention, is
\begin{equation}\tag{LPS5}\label{eq:LPS5}
 \left\|\sum_{q\in\mathscr Q_5} f(r_q^{-1}\,\cdot)\right\|_{L^2(\mathbb S^2)}
             \le2\sqrt5\,\|f\|_{L^2(\mathbb S^2)},
           \qquad\int f=0.
\end{equation}
This is an imported arithmetic theorem, not a finite-dimensional
calculation. The original sources are \cite{LPS86,LPS87}; the
unnormalized norm theorem is explicitly restated in
\cite[Theorem 3(ii), p.\ 8 of the preprint]{Pisier}, while the
quaternionic and Hecke representative conventions are discussed in
\cite[Section 3.2]{PS}. The restated existence theorem alone would
not identify an arbitrary six rotations: the representative convention
in \eqref{eq:LPS5} is essential.

\begin{proposition}[Parity, quotient and full spherical spectrum]
\label{prop:lps-transfer39}
With the external input \eqref{eq:LPS5}, the uniform average over the
six rotations in \eqref{eq:six-gates} has mean-zero norm at most
$\sqrt5/3$ both on $L^2(\mathbb S^2)$ and on $L^2(\SO(3))$.
Its squared gap is at least $4/9$. The conclusion does not assert
the same numerical gap for all of $L^2_0(\SU(2))$.
\end{proposition}
\begin{proof}
An integral quaternion of norm five has one coordinate of magnitude
one, one of magnitude two, and two zero coordinates. There are
$4\cdot3\cdot4=48$ such quaternions. Multiplication by a Lipschitz
unit in $\{\pm1,\pm\mathbf i,\pm\mathbf j,\pm\mathbf k\}$
places the odd coordinate in the real position and makes it positive.
This gives a unique representative with positive odd real part and
even other coordinates, namely one of the six members of $\mathscr Q_5$.
This specifies the parity class; it is not an average over all 48
quaternions or an extra quotient by their unit group.

Normalizing $q$ to $q/\sqrt5$ does not change $r_q$. Quaternionic
conjugation identifies the imaginary unit sphere with
$\SO(3)/\SO(2)$, and $q$ and $-q$ give the same rotation.
For $(1+2\mathbf i)/\sqrt5$ direct multiplication gives
\[
 \mathbf i\mapsto\mathbf i,\quad
 \mathbf j\mapsto(-3\mathbf j+4\mathbf k)/5,\quad
 \mathbf k\mapsto(-4\mathbf j-3\mathbf k)/5.
\]
These are the columns of $R_x$ in \eqref{eq:six-gates}; cyclic
permutation gives $R_y,R_z$. Conjugating a quaternion changes the
sign of its imaginary coordinate and gives its normalized inverse.
Thus the six representatives already include inverse rotations, with
one copy each. Division of the sum in \eqref{eq:LPS5} by six gives
$2\sqrt5/6=\sqrt5/3$, and squaring gives $1-5/9=4/9$.

For the full-group statement, the complexified spherical space is the
orthogonal sum of the harmonic spaces $\mathcal H_\ell$, $\ell\ge0$.
These are precisely all irreducible unitary representations of
$\SO(3)$, one copy of each. The degree-zero space is the constants.
The regular representation on $\SO(3)$ contains the same irreducibles
with multiplicity $2\ell+1$. Convolution acts on each copy by the
same matrix, so its mean-zero operator norm is the supremum of the
same matrix norms on both spaces. No extra invariant subspace is
removed: on the sphere the constant is the only $\SO(3)$-invariant
function. The symmetric alphabet makes the inverse/right versus
left convention immaterial for these norms. Half-integer-spin
representations occur on $\SU(2)$ but not on its $\SO(3)$ quotient;
this argument makes no claim about those extra blocks.
\end{proof}

\paragraph{Source boundary.}
The original LPS publisher records were checked, but their complete
1986/1987 page images were not obtained in this revision. We therefore
do not invent an original theorem number or label the secondary
restatement as an inspected original proof. Equation~\eqref{eq:LPS5}
states exactly the arithmetic line required, and the proposition
proves every subsequent convention change. The effective theorem
retains this explicit external dependency. The main width
classification is conditional on an arbitrary full gap and does not
depend on this numerical example or on an original-page attribution.
